% for landscape posters, use "a4paper, landscape"
\documentclass[a4paper]{article}
\usepackage{../better_poster}
\usepackage{mhchem}
\usepackage{adjustbox}

% ---- fill in from here
% poster size - this will scale the poster to the given size.
% for landscape posters add ", landscape" to the postersize command.
\postersize{a4paper}

\usepackage{polyglossia}
\setdefaultlanguage{czech}

% authors
%\title{A Latex Implementation of Mike Morrison's Better Poster}
%\author{Lana Sinapayen, Someone Else}

% type of poster: [exp]erimental results, [methods], [theory]
% Disclaimer: the original classification had "study" and "intervention" as separate categories. I group them under experimental results.
\newcommand\postertype{exp} % [exp],[methods],[theory]

\begin{document}

% main point of your study
\makefinding{
\textbf{pH} -- tahák\\
\huge https://z-moravec.net/
}


% \makemain{
% }{

% }


% the main text of your poster goes here
\makemain{
    % you can have 1 or 2 columns
    %\raggedcolumns
    %\begin{multicols}{1}

    \section*{Základní vztahy}
    \begin{compactitem}
    	\item Autoprotolýza vody: \ce{H2O + H2O <=> H3O+ + OH-}
    	\item Iontový součin vody: K$_W$ = [H$^+$][OH$^-$] = 10$^{-14}$
    	\item pH = $-$log [\ce{H+}]; pOH = $-$log [\ce{OH-}]
    	\item Pro vodné roztoky platí: pH + pOH = 14
    \end{compactitem}

    \section*{Stupnice pH pro vodné roztoky}

    \begin{compactitem}
    	\item 0--7 kyselé prostředí
    	\item 7 neutrální prostředí
    	\item 7--14 zásadité prostředí
    \end{compactitem}

        \section*{pH kyselin}
        \begin{compactitem}
            \item Silná jednosytná kyselina (HCl): $\textrm{pH = }-\log \textrm{c}_{kys}$
            \item Silná dvojsytná kyselina (\ce{H2SO4}): $\textrm{pH = }-\log (2\times \textrm{c}_{kys})$
            \item Slabá kyselina (HF): $\textrm{pH = }\frac{1}{2}(\textrm{p}K_a - \log \textrm{c}_{kys})$
        \end{compactitem}

\section*{pH zásad}
\begin{compactitem}
	\item Silná jednosytná zásada (NaOH): $\textrm{pH = }14 + \log \textrm{c}_{zas}$
	\item Silná dvojsytná zásada (\ce{Ba(OH)2}): $\textrm{pH = }14 + \log (2\times \textrm{c}_{zas})$
	\item Slabá zásada (\ce{NH3}): $\textrm{pH = 14 } - \frac{1}{2}(\textrm{p}K_b - \log \textrm{c}_{zas})$
\end{compactitem}

\section*{pH solí}
\begin{compactitem}
	\item Soli silné zásady a silné kyseliny (NaCl): \emph{disociací neovlivňují pH}
	\item Soli silné kyseliny a slabé zásady (\ce{NH4Cl}): $\textrm{pH} = 7 - \frac{1}{2}(pK_a + \log c)$
	\item Soli slabé kyseliny a silné zásady (\ce{NaF}): $\textrm{pH} = 7 + \frac{1}{2}(pK_b + \log c)$
	\item Soli slabé kyseliny a slabé zásady (\ce{NH4F}): $\textrm{pH} = 7 + \frac{1}{2}(pK_a - pK_b)$
	\begin{compactitem}
		\item pH solí slabých kyseliny a slabých zásad není závislé na koncentraci soli.
	\end{compactitem}
\end{compactitem}

\section*{pH pufrů}
\begin{compactitem}
	\item Slabá kyselina a její sůl: $\textrm{pH} = pK_a + \log\frac{[A^-]}{[HA]}$
	\item Slabá zásada a její sůl: $\textrm{pH} = 14 - pK_b + \log\frac{[B^+]}{[BOH]}$
\end{compactitem}

    % this determines where your columns will be separated
    %\columnbreak


    %\end{multicols}
}
% If you have extra figures or data to show
\makeextracolumn{
    \textbf{Arrheniova teorie kyseliny a zásad}
    \begin{compactitem}
        \item Kyselina ve vodném roztoku uvolňuje kation \ce{H+}, zásada uvolňuje anion \ce{OH-}.
    \end{compactitem}

	\textbf{Brønstedtova teorie kyseliny a zásad}
	\begin{compactitem}
		\item Kyselina je donorem protonu, zásada jeho akceptorem.
	\end{compactitem}

	\textbf{Jednosytná kyselina}
	\begin{compactitem}
		\item Při disociaci uvolňuje jeden proton:
		\item \ce{HCl -> H+ + Cl-}
	\end{compactitem}

	\textbf{Dvojsytná kyselina}
	\begin{compactitem}
		\item Při disociaci uvolňuje dva protony:
		\item \ce{H2SO4 -> 2 H+ + SO$^{2-}_4$}
	\end{compactitem}

	\textbf{Silná kyselina/zásada}
	\begin{compactitem}
		\item Je v roztoku zcela disociovaná
	\end{compactitem}

	\textbf{Slabá kyselina/zásada}
	\begin{compactitem}
		\item Je v roztoku disociovaná pouze z části
	\end{compactitem}

	\textbf{Sůl}
	\begin{compactitem}
		\item Produkt neutralizace, např. NaCl
	\end{compactitem}

	\textbf{Neutralizace}
	\begin{compactitem}
		\item Reakce kyseliny se zásadou, např.:
		\item \ce{HCl + NaOH -> NaCl + H2O}
	\end{compactitem}

	\textbf{Disociační konstanta}
	\begin{compactitem}
		\item Rovnovážná konstanta disociace kyseliny (K$_a$) nebo zásady (K$_b$). Často se uvádí v logaritmovaném tvaru (pK$_a$ nebo pK$_b$).
		\item \ce{HF <=> H+ + F-}
		\item K$_a = \frac{[H^+][F^-]}{[HF]}$
		\item pK$_a$ = -log K$_a$
	\end{compactitem}

	\textbf{Pufr}
	\begin{compactitem}
		\item Roztok používaný pro stabilizaci pH, je tvořen slabou kyselinou nebo zásadou a její solí
		\item \textit{Acetátový pufr}: směs kyseliny octové a octanu sodného
		\item \textit{Fosfátový pufr}: směs dihydrogenfosforečnanu sodného a hydrogenfosforečnanu sodného
	\end{compactitem}

	\textbf{Stupnice pH}

	\vspace{5mm}

	\adjincludegraphics[width=60mm]{img/PHscalenolang.png}
}

% footer
% generate qr code from https://www.qr-code-generator.com/ and replace qr_code.png
% default: barcode on the left
\makefooter{cc.png}{img/qr-pH.png}

% replace with this like for barcode on the right
%\makealtfooter{images/uni_logo.png}{images/qr-code.png}

\end{document}